% 原创教学示意；近邻数与簇数是不同超参数。
% 左图查询点 (15,13) 的三个最近邻为 (13,15)、(17,11)、(20,16)。
\begin{tikzpicture}[x=1mm,y=1mm,font=\footnotesize]
  \foreach \shift in {0,54}{
    \begin{scope}[xshift=\shift mm]
      \fill[BookPaper] (0,0) rectangle (46,35);
      \foreach \x/\y in {8/8,13/15,17/11,10/19,20/16}{
        \fill[FigureInput] (\x,\y) circle (0.9);
      }
      \foreach \x/\y in {29/23,34/29,38/23,30/30,39/31}{
        \fill[FigureLoss] (\x-0.8,\y-0.8) rectangle (\x+0.8,\y+0.8);
      }
    \end{scope}
  }
  \node[font=\heiti\small] at (23,41) {$k$近邻：$k=3$};
  \node[font=\heiti\small] at (77,41) {K-means：$k=2$};
  \draw[FigureModel,dashed,thick] (15,13) circle (6.7);
  \fill[BookInk] (15,14.2)--(16.2,13)--(15,11.8)--(13.8,13)--cycle;
  \node[align=center] at (23,-7) {保存样本及标签\\每次仅选近邻投票};
  \begin{scope}[xshift=54mm]
    \fill[BookPaper,opacity=0.82] (0,0) rectangle (46,35);
    \foreach \x/\y in {13.6/13.8,34/27.2}{
      \draw[FigureModel,very thick] (\x-2,\y)--(\x+2,\y) (\x,\y-2)--(\x,\y+2);
    }
  \end{scope}
  \node[align=center] at (77,-7) {保存两个中心\\按最近中心分配};
\end{tikzpicture}
